\chapter{Toolchain}\label{ch:toolchain}

The toolchain is built with \texttt{make} and consists of the assembler \texttt{sasm}, the C compiler
\texttt{cc} (as a cross compiler, \texttt{strata-cc}), the reference
emulator \texttt{strata-emu} with its dynamic binary translator, the Verilator-based RTL simulator
\texttt{strata-rtl}, the disk builder \texttt{mkdisk}, the conformance suite in \texttt{conformance/},
the test scripts in \texttt{tools/}, a WebAssembly build of the emulator that runs in a web browser,
and the debuggers: the emulator's debug engine with the text debugger \texttt{strata-dbg} and a web
debugger as its front-ends. \texttt{make} builds \texttt{build/sasm}, \texttt{build/strata-emu},
\texttt{build/strata-rtl} and \texttt{build/strata-dbg}; the
number of RTL cores is fixed at build time by \texttt{NCORES} (default 8), and the \texttt{make} targets
pass the same count to the emulator.

\section{The assembler, \texttt{sasm}}

\begin{shell}
sasm [-o out.stx] [-l] [-g out.map] [-T text-origin] [-I dir] file.s
\end{shell}

\begin{description}[leftmargin=1em,style=nextline]
\item[\texttt{-o}] output image (default: the source name with \texttt{.stx}).
\item[\texttt{-l}] print a listing (section, address, up to four instruction words, source line).
\item[\texttt{-g file}] also write a symbol and line map for the debuggers (Section~\ref{sec:debug}):
  every label with its section and address, every constant, and the address of each source line's
  first byte. The image is the same with or without it.
\item[\texttt{-T addr}] start \texttt{.text} at \texttt{addr} instead of 0. Programs for StrataOS use
  \texttt{-T 0x8000}.
\item[\texttt{-I dir}] add a directory to the \texttt{.include} search path (up to 16).
\end{description}

\texttt{sasm} is a two-pass assembler. It knows every instruction of Part~\ref{part:isa}, including
\insref{svc}, every CSR by name, and the pseudo-instructions of Table~\ref{tab:pseudo}. Its output is an
STX2 image (Section~\ref{sec:stx}). The self-hosted assembler \texttt{as} (Section~\ref{sec:os-as})
accepts the same language and produces byte-identical images.

\subsection{Source syntax}
Each line holds optional labels (\texttt{name:}), then an instruction, directive or assignment.
Comments start with \texttt{;}, \texttt{\#} or \texttt{//} and run to the end of the line. Mnemonics and
directives are case-insensitive. Operands are separated by commas.

\begin{itemize}
\item Integer registers: \reg{x0}--\reg{x31} or ABI names (Table~\ref{tab:abi}), plus \reg{fp} for
  \reg{s0}. Capability registers: \reg{c0}--\reg{c15}. Tags: \reg{tg0}--\reg{tg15}. CSRs: by name
  (Appendix~\ref{app:csrs}) or by number.
\item Memory operands: \texttt{imm(reg)} or \texttt{(reg)}.
\item Symbols start with a letter, \texttt{\_}, \texttt{.} or \texttt{\$}. \texttt{.} alone is the
  current location in the current section.
\item Numbers: decimal, \texttt{0x} hexadecimal, leading-\texttt{0} octal, \texttt{0b} binary; character
  literals \texttt{'a'}, \texttt{'\textbackslash n'}.
\item Expressions, from lowest to highest precedence: \texttt{|}, \texttt{\textasciicircum},
  \texttt{\&}, \texttt{<<} \texttt{>>}, \texttt{+} \texttt{-}, \texttt{*} \texttt{/} \texttt{\%}, unary
  \texttt{-} \texttt{+} \texttt{\textasciitilde}, parentheses. Arithmetic is 64-bit.
\end{itemize}

Immediates are range-checked: signed 16 bits for arithmetic, loads, stores and \insn{jalr}; unsigned
16 bits for logical immediates; 0--63 for shifts. Branch and jump targets must be in the same section
and in range.

\subsection{Sections}
Each section has its own address space starting at 0:

\begin{center}\small
\begin{tabular}{@{}l l l@{}}
\toprule
\textbf{Directive} & \textbf{Assembles into} & \textbf{Labels are}\\\midrule
\texttt{.text} & instruction Local Store & I-LS addresses (from \texttt{-T})\\
\texttt{.data} & data Local Store & D-LS addresses\\
\texttt{.far} & T3 system memory & offsets in T3\\
\texttt{.cluster} & T2 cluster store & offsets in T2\\
\bottomrule
\end{tabular}
\end{center}

\subsection{Directives}

\begin{table}[htbp]
\centering\small
\caption{Assembler directives}\label{tab:directives}
\begin{tabularx}{\linewidth}{@{}>{\ttfamily}l X@{}}
\toprule
\textrm{\bfseries Directive} & \textbf{Effect}\\\midrule
.text .data .far .cluster & select a section\\
.org expr & advance the location to \textit{expr} (never backwards), filling with zeros\\
.align n & pad with zeros to a multiple of $n$ bytes ($n$ a power of two)\\
.byte .half .word .dword e, \ldots & emit 1-, 2-, 4- or 8-byte little-endian values\\
.ascii "s" & emit a string; escapes \texttt{\textbackslash n \textbackslash t \textbackslash r \textbackslash 0 \textbackslash x}\textit{hh}\\
.asciz "s", .string "s" & emit a string followed by a zero byte\\
.space n[, fill], .zero n & emit $n$ bytes of \textit{fill} (default 0)\\
.equ name, expr; .set; name = expr & define a constant (no forward references)\\
.include "file" & include a file (searched next to the including file, then the current directory, then \texttt{-I} dirs)\\
.global, .globl & accepted and ignored\\
\bottomrule
\end{tabularx}
\end{table}

\section{The STX2 image format}\label{sec:stx}

An image holds the initial contents of each section. Every part of the file is 8-byte aligned so that
software (the boot ROM, the OS loader) can move it with transfers.

\begin{center}\small
\begin{tabular}{@{}r l l@{}}
\toprule
\textbf{Offset} & \textbf{Type} & \textbf{Field}\\\midrule
0 & char[4] & magic \texttt{"STX2"}\\
4 & u32 & number of sections\\
\multicolumn{3}{@{}l}{\textit{then, for each section:}}\\
+0 & u32 & kind: 1 I-LS, 2 D-LS, 3 T3, 4 T2\\
+4 & u32 & base: lowest address used, rounded down to a multiple of 8\\
+8 & u32 & size in bytes, a multiple of 8\\
+12 & u32 & reserved, zero\\
+16 & u8[size] & contents of $[\textit{base}, \textit{base}+\textit{size})$\\
\bottomrule
\end{tabular}
\end{center}

Empty sections are omitted. When an image is loaded by a simulator, the I-LS and D-LS contents go to
every core, and the T2 and T3 contents to the shared memories.

\section{The C compiler, \texttt{cc}}\label{sec:cc}

\texttt{cc} (\texttt{lang/cc/}) compiles C99 to StrataOS program images. It is one pass over one source
file: code is emitted while the source is parsed, the image is written directly, and there is neither
an assembler nor a linker behind it. Its source, \texttt{lang/cc/cc.c}, is written in the C it compiles
and is used twice: the host's C compiler builds it, with a small layer standing in for StrataOS
(\texttt{lang/cc/host/}), into the cross compiler \texttt{build/strata-cc}; and compiled by itself it is
the \texttt{cc} on the StrataOS disk (Section~\ref{sec:os-as}). Both write the same images, byte for
byte.

\begin{shell}
strata-cc [-o out] [-p] [-h kib] file.c
\end{shell}

\begin{description}[leftmargin=1em,style=nextline]
\item[\texttt{-o}] the image (default: the source's name without \texttt{.c}).
\item[\texttt{-p}] the paged code model: every function is placed in a 4~KiB code unit of the image's
  T3 section, calls between functions are \texttt{xcall} and returns \texttt{xret}, and StrataOS pages
  the units into the instruction Local Store (Section~\ref{sec:os-paging}). Without it the program and
  the library share the 32~KiB of I-LS above \hex{8000}. An image holds at most 128~KiB; the units of a
  larger program go to a second file, \textit{out}\texttt{.u}, which the start-up code loads.
\item[\texttt{-h}] the size of the heap in KiB (default: what the data Local Store has left).
\end{description}

\paragraph{The machine as a C target.} \texttt{int} is 32 bits, \texttt{long}, \texttt{long long} and
pointers are 64. A pointer is an address in the core's data Local Store: globals, string literals, the
heap and the stack (8~KiB) share its 40~KiB of program data. System memory is not addressed: it is
reached by transfers, through the library (\texttt{far\_read}, \texttt{far\_write}) or the intrinsics.
Loads and stores must be aligned, so the compiler never emits a wide access to a narrowly aligned
object. Arguments are passed in \texttt{a0}--\texttt{a7} and then on the stack, results returned in
\texttt{a0}; \texttt{gp} holds \hex{10000}, which puts every static object and every library routine
one instruction away. There is no floating-point hardware: \texttt{float} and \texttt{double} are
IEEE~754 in software, in the library.

\paragraph{The library.} \texttt{lang/cc/lib/libc.s}, in assembly, is assembled once at fixed addresses
and copied into every image: the start-up code, files and the console over StrataOS messages,
\texttt{printf}, the heap, strings, the software floating point. \texttt{cc.lib} holds it with the
headers. Functions a program may never call are kept out of it: \texttt{qsort}, \texttt{bsearch} and
the functions of \texttt{<math.h>} are C sources in \texttt{cc.lib} (\texttt{lang/cc/lib/src/}), and
when a program has been read and still calls one that nothing defines, the compiler compiles that
source as if it followed the program.

\paragraph{The machine from C.} The instructions only \IS has are intrinsics, one instruction each
(\texttt{<strata.h>}): the transfers \texttt{\_\_tget}, \texttt{\_\_tput}, \texttt{\_\_twait}, the far
atomics, the capability instructions, CSR access, \texttt{\_\_svc}. Anything else is
\texttt{asm("...")}: assembly in the \texttt{sasm} syntax, encoded in place, in which
\texttt{\{name\}} stands for a C object, constant or function; at file scope its labels are functions.
The table of instructions is not written into the compiler: the build derives it from \texttt{sasm},
by assembling each instruction and seeing which bits each operand changes.

\paragraph{Two compilers.} A machine with no C compiler cannot compile \texttt{cc.c}, so there is a
second compiler of the same language, in \IS assembly (\texttt{lang/cc/cc.s}, the program \texttt{cc0}
of the disk). It compiles \texttt{cc.c}; the result compiles \texttt{cc.c} again; and what that
produces, compiling \texttt{cc.c} once more, reproduces itself and equals the cross-built \texttt{cc}.
\texttt{tools/test-cc-boot.sh} runs this inside StrataOS and compares the files there.

\paragraph{Tests.} \texttt{lang/cc/tests/} holds 521 C programs from the test suites of two other
compilers (scc and lacc), ordered from integer expressions to floating point. Each is compiled and run
inside StrataOS, and its exit status and output must be those of the reference; both compilers pass
all of them in both code models (\texttt{make test-cc}, \texttt{test-cc0}; \texttt{test-cc-host}
compiles with \texttt{strata-cc} and only runs the programs in the emulator). \texttt{make test} runs a
subset, the bootstrap, and \texttt{tools/test-cc.sh}, which compiles and runs the examples of
\texttt{lang/cc/examples/} at the StrataOS shell.

\section{The reference emulator, \texttt{strata-emu}}\label{sec:emu}

\begin{shell}
strata-emu [-c ncores] [-t trace.txt] [-m max_steps] [-d disk.img] [-i input.txt]
           [-b steps] [-T real|virtual] [-g none] [-F frame.ppm] [-z scale]
           [-r WxH] [-S entries] [-X entries] [-L D[,I]] [-H] [-P pointer.txt] [-N mode] [-J|-j]
           [--debug where [--map file] [--os dir] [--www dir] [--history MiB]] image.stx
\end{shell}

The emulator executes the sequential semantics $\mathcal{S}$ (Chapter~\ref{ch:model}): every memory
transfer completes at issue, and the only transfers that stay in flight are device reads that have
nothing to deliver yet (console input, the break port, the timer's WAIT and the framebuffer
controller's VSYNC; Section~\ref{sec:devxfer}), at most four per core. \insn{twait} and \insn{tpoll}
see only those reads, and \csr{cycle} equals \csr{instret}. Cores take turns one instruction at a time
(a \term{round} is every live core taking one step); a core blocked in \insn{cgive}, \insn{csend},
\insn{ctake}, \insn{crecv}, \insn{tany}, on a busy tag or on a Local Store range a pending device read
is filling yields its turn. The emulator reports each halt (\texttt{core N: halt a0=\ldots
instret=\ldots}) and each unhandled fault on standard error. When every live core is blocked it
detects deadlock, unless a pending read can still complete: a WAIT or VSYNC whose time will come, or
console input or Ctrl-C from an interactive terminal or the window.

\begin{description}[leftmargin=1em,style=nextline]
\item[\texttt{-c n}] number of cores, 1--15 (default 4). Fifteen cores and the block device are the
  sixteen inbox senders that the masks of \insn{cpoll} and \insn{tany} can name.
\item[\texttt{-t file}] instruction trace, one line \texttt{core pc word} per executed instruction.
  Implies interpretation.
\item[\texttt{-m n}] stop with a timeout after $n$ instructions in all; by default there is no limit.
\item[\texttt{-d file}] disk image for block device 0.
\item[\texttt{-i file}] console input (default standard input). When standard input is a terminal it
  is put in character mode: keys are delivered as typed, unechoed, and Ctrl-C goes to the console break
  port. At the end of the file the input port returns \hex{04} forever.
\item[\texttt{-P file}] scripted pointer input (Section~\ref{sec:pointer}), one event per line:
  \textit{when} \texttt{move} $x$ $y$ $|$ \texttt{down} $b$ $|$ \texttt{up} $b$ $|$ \texttt{wheel} $dy$
  [$dx$] $|$ \texttt{enter} $|$ \texttt{leave}, where $b$ is \texttt{left}, \texttt{right} or
  \texttt{middle}. \textit{when} is an instruction count (on \texttt{strata-rtl}, a cycle count), or
  \texttt{k}$n$: once the program has taken $n$ bytes of console input. Lines happen in order.
\item[\texttt{-N mode}] the network interface's link (Section~\ref{sec:netif}; \texttt{common/net.h}):
  \texttt{none} (default), \texttt{loop} (packets come straight back; deterministic),
  \texttt{peer} (a built-in host at every address but the sender's own and 192.0.2.0/24: it
  answers ICMP echo, echoes UDP datagrams sent to port~7 and answers DNS for its names
  \texttt{peer.test}, \texttt{two.test}, \texttt{www.peer.test}; other names are looked up for real,
  or with \texttt{peer,nodns} do not exist), \texttt{nat}[\texttt{,fwd=}\textit{hostport}\texttt{:}%
  \textit{guestport}]\ldots{} (the real network through a user-mode NAT, \texttt{common/netnat.h}:
  the guest is 10.0.2.15, the host 10.0.2.2, DNS 10.0.2.3; UDP and TCP through host sockets, TCP
  terminated on the host and relayed with retransmission toward the guest, ICMP echo through an
  unprivileged ICMP socket where the host has one; each \texttt{fwd} listens on
  127.0.0.1:\textit{hostport} for the guest's \textit{guestport}, and a busy port is reported and
  skipped), or \texttt{udp:}\textit{lport}\texttt{:}\textit{host}\texttt{:}%
  \textit{rport}: one packet per host UDP datagram, sent from \textit{lport} to
  \textit{host}:\textit{rport} and received only from there. Two simulators with mirrored ports are
  two machines on one link. None needs privileges. \texttt{make boot} passes
  \texttt{-N nat,fwd=8080:80,fwd=8000:8000}. \texttt{tools/netrelay} is the same NAT behind a
  WebSocket server, for the browser's play page (one IPv4 packet per binary message).
\item[\texttt{-b n}] press Ctrl-C once, after $n$ instructions: it completes the oldest pending break
  read, or is remembered for the next one. Used by \texttt{tools/test-break.sh}.
\item[\texttt{-T real|virtual}] the timer's time base (Section~\ref{sec:timer}). \texttt{real}
  (default): host wall-clock microseconds since start. \texttt{virtual}: rounds / 100, a nominal
  100~MHz machine; when every core waits only for the timer, virtual time skips ahead to the next
  deadline. Runs whose output is compared use \texttt{-T virtual}.
\item[\texttt{-g none}] never open a window. PRESENT is still accepted and still records the frame
  for \texttt{-F}. (\texttt{-g window} is the default.)
\item[\texttt{-F file}] at exit, write the frame as of the last PRESENT as a binary PPM. Writes into the
  framebuffer after that PRESENT do not change the file.
\item[\texttt{-z n}] window scale, 1--8 (default 1).
\item[\texttt{-r W}x\texttt{H}] framebuffer geometry (default 1280\,x\,800; width a multiple of 8, from
  320\,x\,200 to 2560\,x\,1600). The platform table reports it (Section~\ref{sec:stpd}).
\item[\texttt{-S n}, \texttt{-X n}] the capability save table's entries per core ($\textit{CTAB}$,
  64--256, default 256) and the overlay map's ($\textit{NXMAP}$, 4--16, default 16).
\item[\texttt{-L} \textit{D}[\texttt{,}\textit{I}]] each core's D-LS and I-LS in KiB
  (Section~\ref{sec:lssize}): powers of two, 4--64 (default 64); \textit{I} defaults to \textit{D}. An
  image whose Local Store sections do not fit is refused. The platform table reports all three
  parameters, and an RTL build has them as parameters (\texttt{make NXMAP=}, \texttt{DLS\_KIB=},
  \texttt{ILS\_KIB=}); runs that compare the two pass the same values.
\item[\texttt{-H}] when every core has halted, keep the window open, showing the final memory contents,
  until it is closed.
\item[\texttt{-J}, \texttt{-j}] run with the JIT, or interpret only. The default is the JIT, except
  with \texttt{-t} or \texttt{-T virtual}, which keep the exact one-instruction interleaving.
\item[\texttt{-{}-debug} \textit{where}] run under the debug engine (Section~\ref{sec:debug}), stopped in
  front of the first instruction, controlled through the debug protocol: \textit{where} is
  \texttt{stdio}, a TCP port (\texttt{PORT} or \texttt{HOST:PORT}: JSON lines, a WebSocket and the web
  debugger's pages on one port) or \texttt{unix:}\textit{path}. The companion options:
  \texttt{-{}-map} \textit{file} loads a symbol map (\textit{image}\texttt{.map} is loaded by itself);
  \texttt{-{}-os} \textit{dir} turns StrataOS awareness on (\textit{dir}\texttt{/kernel.map},
  \textit{dir}\texttt{/prog/}\textit{name}\texttt{.map}); \texttt{-{}-www} \textit{dir} is what the
  port serves to a browser (default \texttt{site/}); \texttt{-{}-history} \textit{MiB} bounds time
  travel's memory (default 512).
\end{description}
The exit status is 0 when every core halted (or the window was closed), 1 if any core faulted, on
deadlock or on timeout, and 2 for a usage error.

\subsection{The graphics window}\label{sec:window}

Both simulators open a window (SDL2) the first time software writes into the framebuffer window of
system memory (a \insn{tput}, a \insn{tamo} or a block-device read landing in
\hex{3000000}--\hex{3FFFFFF}) or writes PRESENT, whichever comes first; before that they are plain
terminal programs. The window scans out the framebuffer memory live (height $\times$ stride bytes,
Section~\ref{sec:fbctl}): it is repainted from the current memory contents at most every 20~ms of host
wall-clock time (about 30~Hz for \texttt{strata-rtl}, which reads back only the range written since the
last repaint), only when something was written since then, on expose, and at once on PRESENT. Closing
the window ends the simulation. A build without SDL2 (\texttt{-DNO\_SDL}) has no window and still
records frames for \texttt{-F}.

While the window has focus, keys are fed into the console input stream as a terminal sends them:
printable characters as themselves; Enter \hex{0A}, Backspace \hex{7F}, Tab \hex{09}, Escape \hex{1B};
the arrows \texttt{ESC [ A}/\texttt{B}/\texttt{C}/\texttt{D} (up, down, right, left); Home
\texttt{ESC [ H}; End \texttt{ESC [ F}; Delete \texttt{ESC [ 3 \textasciitilde}; Page Up and Page Down
\texttt{ESC [ 5 \textasciitilde} and \texttt{ESC [ 6 \textasciitilde}; Ctrl+letter as 1--26, except
Ctrl+C, which is the break key (Section~\ref{sec:console}); Alt+key in the xterm modified-key forms:
Alt with a character as \texttt{ESC [ 27 ; 3 ;} \textit{code} \texttt{\textasciitilde} (the character's
code in decimal, Shift applied: Alt+Shift+Q is 81; Alt+Enter is 13), Alt with an arrow, Home or End as
\texttt{ESC [ 1 ; 3 A}/\texttt{B}/\texttt{C}/\texttt{D}/\texttt{H}/\texttt{F}, and Alt with Delete,
Page Up or Page Down as \texttt{ESC [} \textit{n} \texttt{; 3 \textasciitilde}. Not ESC followed by the key
(``meta sends escape''): those are also the bytes of Esc and then the key, which a reader of a byte
stream could only tell apart by time. Nobody types the modified forms by hand, so the window server
takes only them as its Alt bindings, and Esc and then a key is just that.
The terminal remains an input source at the same time.

\subsection{The JIT}\label{sec:jit}

\texttt{emu/jit/} is a dynamic binary translator with an arm64 and an x86-64 backend, selected at
compile time; on other hosts the emulator only interprets. It translates the plain integer instruction
set (ALU, immediates, D-LS loads and stores, branches, \insn{jal}, \insn{jalr}) into host code, a
superblock at a time (up to 128 instructions, running on past conditional branches, which leave
through side exits), and chains blocks through a direct-mapped table with one entry per I-LS word.
Every other instruction, and every instruction that traps, is left to the interpreter: translated code
stops in front of it with the core's state exactly as the interpreter would have it, so traps are
precise and \csr{instret} is exact. A \insn{tgetx} invalidates every translation that read the
overwritten I-LS range. With the JIT a core runs a quantum of up to 8192 instructions per turn, and
ends its turn early at an instruction that meets other cores (a device or core-control transfer, a far
atomic, \insn{tany}, \insn{cgive}, \insn{ctake}, \insn{cpoll}, \insn{csend}, \insn{crecv}); a core
with a pending device read or a pending stop is interpreted. Results of programs that do not depend
on timing (Section~\ref{sec:determinism}) are identical with and without the JIT; \texttt{make
test-jit} checks this on every example and on the Forth regression test.

\section{The debuggers}\label{sec:debug}

The emulator has a debug engine (\texttt{emu/debug.c}) that runs a machine under control, and two
front-ends: \texttt{strata-dbg}, a text debugger in the manner of gdb, and a web debugger
(\texttt{site/debug/}). Nothing in the architecture changes for them: breakpoints and watchpoints live
in the emulator, and without \texttt{-{}-debug} it runs exactly as before. The protocol between them,
JSON objects one per line, is documented in \texttt{doc/debug-protocol.md}.

\subsection{The engine}
The engine runs the reference interleaving: round after round, each live core executes one instruction
or finds itself blocked. A \term{tick} is one core's turn in a round; ticks are counted from power-on
and are the debugger's clock. Before a core's turn the engine decodes the instruction at its \reg{pc}
(breakpoints stop there, in front of the instruction); after it, it compares the core with what it was
(transfers issued, capabilities made and given, bytes written) and looks at every core's pending
device reads (completions). The emulator reports to it only what cannot be seen from outside: writes to
far memory and to the disk, traps and console output.

\begin{description}[leftmargin=1em,style=nextline]
\item[Control] continue, pause; \texttt{stepi} one core (the others take their turns in the same
  rounds), or the thread of a StrataOS process across switches and cores; \texttt{next} over a call,
  \texttt{finish}, \texttt{until} a location, whole rounds. A plain continue with no breakpoint set,
  on a run with the JIT, runs at full speed and stops between quanta.
\item[Breakpoints] at a \reg{pc} (on any core or one, with an ignore count or a condition on a register,
  under StrataOS in one program only), at a trap cause, \insn{svc}, \insn{eret}, a halt, a capability
  operation (\insn{cgive}, \insn{ctake}, \insn{csend}, \insn{crecv}, \insn{cdom}, \insn{crevoke},
  \ldots), or when a StrataOS process starts to run.
\item[Watchpoints] on D-LS or I-LS bytes (written by a store or by a transfer at its completion, read by
  a load or a \insn{tput}), on far ranges of T2, T3 or device space (touched by transfers, atomics and
  the block device), on a capability register, on a tag's issue or completion.
\item[State] registers and CSRs; the capability registers decoded (space, base, length, permissions,
  revocable and generation, live or dead), the save table, the inboxes, the overlay map, the revocation
  units; the sixteen tags of every core (the last transfer on each: operation, Local Store range, far
  space and offset, length, issue and completion round); the pending device reads and what each waits
  for; the block device's queues; why a core is stalled (which tag, which hazard, which inbox); memory
  of any space; disassembly with labels and source lines; a backtrace by heuristics.
\item[Event log] transfers issued and completed, traps, \insn{eret}, gifts and messages, revocations,
  stalls and halts, for the timelines of the front-ends.
\item[Provenance] every capability the machine makes is given an id, which moves keep, and the engine
  records its parent and the operation that made it (split, share, restrict, give, mint, revoke,
  \ldots), so a front-end can show where any capability came from.
\item[Time travel] on a deterministic run (\texttt{-T virtual}, no network) the engine keeps
  snapshots of the machine (the pages of memory that changed since the last one) and a log of what came
  from outside (console input, Ctrl-C, the debugger's own writes). Going back is restoring the nearest
  earlier snapshot and executing forward again: to any tick, back one step of a core, back to the last
  breakpoint or watchpoint hit, back to the issue of a tag. The memory it uses is bounded; old snapshots
  are thinned out. A fast continue catches up without recording (with the JIT under \texttt{-J}) and
  recording starts again where it stops, so a boot need not be recorded to debug what follows it.
\item[StrataOS] with \texttt{-{}-os}, the kernel's shared table and \texttt{proc}'s process table are
  read through the maps (no layout is hard-coded): processes, the process each core runs, the program
  map for user code, the kernel's code units by name through core 0's overlay map.
\end{description}

Symbols come from the maps \texttt{sasm -g} writes (the Makefile writes one for every image, every
program on the StrataOS disk, the kernel and the boot ROM); Strata Forth's code is also named from its
dictionary. In the emulator a memory transfer completes the instant it issues (Chapter~\ref{ch:model}),
so only device reads are ever seen in flight.

\subsection{\texttt{strata-dbg}}

\begin{shell}
strata-dbg [--console pty|file|-] [--tui] [-x script] [--batch] [emulator options] image.stx
strata-dbg [options] --connect host:port | unix:path
\end{shell}

\texttt{strata-dbg} starts \texttt{strata-emu -{}-debug stdio} with the emulator's options, or connects
to an engine that runs already. The guest's console never mixes with the debugger: by default it is a
new pseudo-terminal, to be opened in another window (what is typed there is the guest's input);
\texttt{-{}-console} \textit{file} sends it to a file or another terminal, and \texttt{-} prints it
inline, each line marked (the default for scripts). Its commands:

\begin{center}\small
\begin{tabularx}{\linewidth}{@{}>{\ttfamily}l X@{}}
\toprule
\textrm{\bfseries Command} & \textbf{Effect}\\\midrule
continue, stepi [n], next, finish & run; step the current core; step over a call; until the function returns\\
until \textit{loc}, step-round [n] & run to a location; run whole rounds\\
core \textit{n}, follow \textit{pid} & make a core current; let \texttt{stepi} follow a process\\
reverse-stepi [n], reverse-continue & back one step of the current core; back to the last breakpoint hit\\
reverse-step-round [n] & back to the start of the round, or of earlier ones\\
reverse-tag tg\textit{n}, goto \textit{tick} & back to a tag's last issue; to any tick of the run\\
break, tbreak \textit{loc} & \textit{loc}: an address, \texttt{label}, \texttt{label+n}, \texttt{file.s:line}, \texttt{prog:label}; then optionally \texttt{core} \textit{n} and \texttt{if} \textit{reg op value}\\
break trap, svc, eret, halt, capop, process & the other kinds of breakpoint\\
watch, rwatch, awatch & Local Store bytes, far ranges (\texttt{t2}, \texttt{t3}, \texttt{dev}), \texttt{cap c}\textit{n}, \texttt{tag tg}\textit{n}\\
delete, enable, disable & manage them (\texttt{info breakpoints} lists them)\\
info registers, caps, ctab, inbox, xmap & a core's state\\
info tags [all], pending, cores & transfers, device reads, every core and why it is stalled\\
info processes, granules, history & StrataOS processes; revocation units; the recorded run\\
x/\textit{nf} [\textit{space}:]\textit{addr} & memory: units \texttt{b h w g}, \texttt{s}, \texttt{i}; spaces \texttt{dls ils t2 t3 disk}\\
disassemble, list, backtrace & code, source, the call chain\\
log [n], prov c\textit{n} & the last events; a capability's history\\
set \textit{reg value}, input \textit{text}, ctrl-c & change a register; console input; the break key\\
tui & the cores side by side: disassembly, registers, capabilities and tags\\
\bottomrule
\end{tabularx}
\end{center}

\subsection{The web debugger}
\texttt{site/debug/} shows the cores side by side (disassembly with the current line and breakpoints,
registers, capabilities, stall reasons), a transfer timeline (a row per core: transfers as they issue,
device reads in flight from issue to completion, stalls, gifts between cores, traps), a Local Store map,
the provenance graph of the capabilities held, breakpoint and watchpoint management, the console and a
slider over the recorded run. In the page it runs the WebAssembly emulator in a worker with
\texttt{-{}-debug}; served by \texttt{strata-emu -{}-debug} \textit{port} it drives the desktop engine
over a WebSocket on the same port (\texttt{make debug-web}). On wide screens the transfers are docked
on the right, beside the cores, with a list of the newest ones. Its first program is StrataOS: the play
page's disk booted with \texttt{-{}-os}, and a Run menu that types a program's name at the shell and can
stop at the process's first instruction (\texttt{make debug-web-os} on the desktop). In the page the
console is a terminal and the framebuffer opens when the guest first draws; keys and the pointer go through
the engine (\texttt{input}, \texttt{pointer}), which logs them for time travel.

\section{The RTL simulator, \texttt{strata-rtl}}

\begin{shell}
strata-rtl [-t trace.txt] [-m max_cycles] [-d disk.img] [-i input.txt] [-b cycles]
           [-T real|virtual] [-g none] [-F frame.ppm] [-z scale] [-r WxH] [-H] [-N mode] image.stx
\end{shell}

The SystemVerilog implementation (Chapter~\ref{ch:impl}) compiled with Verilator 5 and the harness
\texttt{sim/main.cpp}. The number of cores is fixed at build time (\texttt{make NCORES=}$n$, default
8), and so are the overlay map's entries and the Local Stores' sizes (\texttt{make NXMAP=}$n$,
\texttt{DLS\_KIB=}$d$ \texttt{ILS\_KIB=}$i$: the emulator's \texttt{-X} and \texttt{-L}). It accepts the emulator's options with these differences: \texttt{-m} and \texttt{-b} count clock
cycles; \texttt{-T virtual} makes the timer count cycles / 100 (\texttt{-T real}, the default, reads
the host clock through a DPI call); there is no \texttt{-c} and no JIT. A console read with no byte
available is retried by the RTL without stopping the simulation. Halt lines carry the core's cycle
and stall counts in addition to the emulator's fields, and traces have the same format.

\section{Cross-checking}\label{sec:compare}

\texttt{tools/compare.sh image.stx \ldots} runs each image on both machines with \texttt{-T virtual -g
none} and requires that they agree on: the exit status; the console output, ignoring lines that begin
with \texttt{@} (by convention those carry cycle counts or times); each core's halt or fault line,
without the RTL's cycle and stall counts; and each core's instruction trace. Neither machine may time
out or deadlock. Because the two machines have entirely different timing, agreement is evidence both
that the RTL implements the semantics and that the programs satisfy the determinism conditions of
Section~\ref{sec:determinism}.

\texttt{tools/compare-boot.sh} boots StrataOS from the disk image on both machines with \texttt{-T
virtual -g none}, feeds the scripted session \texttt{os/demo.txt} to the console, and compares the
console output and, for every core, the instruction trace of the boot ROM up to the jump into the
kernel. StrataOS preempts processes on timer ticks, and the two machines' clocks advance at different
rates per instruction, so the moments of context switches differ; the lines that report timing (the
figures of \texttt{ps}, \texttt{cores}, \texttt{uptime} and \texttt{sched}) are normalised, and
everything the programs compute must be identical.

Further scripts:
\begin{description}[leftmargin=1em,style=nextline]
\item[\texttt{tools/test-break.sh}] boots StrataOS on both machines, starts \texttt{spin}, presses
  Ctrl-C partway with \texttt{-b}, then runs another program; both machines must print the same output,
  show the program killed, and show that the shell survived.
\item[\texttt{tools/test-preempt.sh}] two \texttt{spin} processes pinned to one worker core both make
  progress, and \texttt{ps} shows context switches and bytes moved for each.
\item[\texttt{tools/test-boot-jit.sh}] the scripted session three times with \texttt{-J -T real}: coarse
  interleaving and host-time preemption; each run must end in a halt, without a fault, with the
  programs' results present.
\item[\texttt{tools/test-fb.sh}] runs \texttt{gfx} on both machines with \texttt{-g none -F} and
  compares the two presented frames byte for byte.
\item[\texttt{tools/test-jit.sh}] every example and the Forth regression input with \texttt{-J} and
  \texttt{-j}: same output, exit status, halt lines and frames.
\item[\texttt{tools/test-sh.sh}] a scripted session of the interactive shell (line editing, history,
  completion, jobs) compared with \texttt{tools/test-sh.out}.
\item[\texttt{tools/test-as.sh}] the self-hosted assembler against \texttt{sasm}
  (Section~\ref{sec:os-as}).
\end{description}

\section{The conformance suite}\label{sec:conformance}

\texttt{conformance/} holds the architectural conformance suite. Each test in
\texttt{conformance/tests/} is a bare-metal program that exercises one area of the architecture and
checks its own results against values derived from this manual: the ALU, loads and stores, and branches
(generated by \texttt{conformance/gen/*.py}, which compute the expected values independently of every
implementation), transfers and hazards, pending-read ordering and \cause{pend} traps, capabilities and
their alignment, revocation by address and endpoints, multicore gifts, core windows and stops, traps,
\insn{svc} and the user-mode inbox rule, code paging, reserved fields and CSRs, halt draining, and every
device: console, break and release, platform table, block device, timer and framebuffer controller,
and device transfers spanning several ports.

A test prints a \term{signature} (lines \texttt{sig} \textit{hex}) and a verdict, \texttt{pass} or
\texttt{fail} with one \texttt{bad} line per failed check. \texttt{conformance/run.sh} runs every test
on three machines at once: the emulator interpreting (\texttt{-j}), the emulator with the JIT
(\texttt{-J}) and the RTL, all with \texttt{-T virtual -g none}, and prints a table of tests by machine.
A cell passes when the machine ran the test to completion, printed \texttt{pass}, and, for the JIT and
the RTL, printed exactly the interpreter's signature; otherwise the cell names the reason
(\texttt{FAIL}~$n$, \texttt{SIG}, \texttt{NOVERD}, \texttt{TIMEOUT}, \texttt{FAULT}, \texttt{DEADLOCK}).
A test may give extra simulator arguments (such as console input) on a line \texttt{; args:} in its
header. Cells listed in \texttt{conformance/xfail.txt} (\textit{test machine reason}) are reported as
\texttt{xfail} and do not fail the run; a listed cell that passes is reported as \texttt{XPASS}.

\texttt{conformance/fuzz/fuzz.py} is a three-way differential fuzzer. Each seed generates a random
valid program: every ALU function on random and edge values, loads and stores of every width (some out
of range or misaligned), transfers between the D-LS and system memory with random tags, lengths and
overlaps followed at once by accesses to the same bytes, far atomics, branches and calls, and a trap
handler that folds every cause, \csr{epc} and \csr{tval} into a register. The program prints its
registers, a hash of its memory and its trap count; the three machines must print the same bytes and
halt with the same \reg{a0}. A failing seed's program is kept for reproduction.

\section{The emulator in a web browser}\label{sec:web}

\texttt{make wasm} compiles the emulator with Emscripten (interpreter only: the JIT is not built) into
\texttt{build/web/}, and \texttt{make web-site} assembles the project's site with it under
\texttt{play/}. The page drives the emulator from its animation loop, running rounds for about 10~ms
per frame, so the browser never blocks; when every core is blocked the frame ends at once. Console
output and input go through an in-page terminal (xterm.js) with the keyboard mapping of
Section~\ref{sec:window}, the framebuffer is painted onto a canvas from live memory (only the rows
written since the last paint), and a status line shows the instruction rate, the cores and the virtual
time. Presets boot StrataOS on 12 cores, run \texttt{gfx}, \texttt{hello} or \texttt{mandel} on 8
cores, or start Strata Forth on 8 cores with a 640\,x\,400 screen and its demos
(Chapter~\ref{ch:forth}); Reset and Ctrl-C buttons complete the controls.

\section{Build targets}

\begin{center}\small
\begin{tabularx}{\linewidth}{@{}>{\ttfamily}l X@{}}
\toprule
\textrm{\bfseries Target} & \textbf{Effect}\\\midrule
make & build \texttt{build/sasm}, \texttt{build/strata-emu}, \texttt{build/strata-rtl}, \texttt{build/strata-dbg}\\
make examples & assemble every example\\
make test & every example on both machines (\texttt{compare.sh}); StrataOS boot compared (\texttt{compare-boot.sh}); Ctrl-C (\texttt{test-break.sh}); preemption (\texttt{test-preempt.sh}); StrataOS with the JIT (\texttt{test-boot-jit.sh}); framebuffer frames (\texttt{test-fb.sh}); JIT against interpreter (\texttt{test-jit.sh}); the debuggers (\texttt{test-dbg.sh}); the conformance suite\\
make conformance & the conformance suite on the interpreter, the JIT and the RTL\\
make fuzz & the differential fuzzer (\texttt{SEEDS=}$n$, \texttt{FIRST=}\textit{seed})\\
make test-jit & the JIT against the interpreter\\
make bench & wall time of the interpreter and the JIT on \texttt{mandel}, \texttt{gfx} and the Forth graphics demo\\
make run, make run-emu & show every example's output on the RTL, or the emulator\\
make run-NAME & run one example on the RTL\\
make show-NAME, show-rtl-NAME & run one example on the emulator (real-time timer) or the RTL and keep its window open after it halts, e.g.\ \texttt{make show-gfx}\\
make lint & Verilator lint of the RTL\\
make os & build the boot ROM and the disk image\\
make boot & boot StrataOS on the emulator (12 cores), interactively\\
make boot-demo & boot StrataOS on the emulator with the scripted session \texttt{os/demo.txt}\\
make boot-rtl, boot-rtl-demo & the same on the RTL simulator (\texttt{NCORES} cores)\\
make ws-demo & the window server's renderers as a bare-metal demo in a window (Chapter~\ref{ch:os})\\
make ws-shots, ws-check, ws-perf & window-server screenshots; damage-driven frames compared with full redraws; instruction counts\\
make forth, forth-rtl & the Strata Forth REPL on the emulator or the RTL\\
make forth-demo-NAME, forth-demos & one Forth demo (\texttt{tiers}, \texttt{multicore}, \texttt{gfx}), or all three, on the emulator\\
make test-forth & the Forth regression test on both machines\\
make wasm, web-site & the WebAssembly emulator; the site with the emulator and the source browser\\
make dbg & \texttt{strata-dbg} and the emulator\\
make debug-NAME, debug-os & one example, or StrataOS (aware of its processes), under \texttt{strata-dbg}\\
make debug-web & the debug engine on port 4444 with the web debugger at \texttt{/debug/}\\
make debug-web-os & the same with StrataOS booting\\
make test-dbg & the debuggers driven by scripts (\texttt{tools/test-dbg.sh}; part of \texttt{make test})\\
make manual & this manual\\
make clean & remove \texttt{build/}\\
\bottomrule
\end{tabularx}
\end{center}
